\subsection{二つの数体曲線を閉じ、共鳴面全体を排除する}
\begin{proposition}
非退化正top・対称secant負top・正Pfaffianの主chartにおいて、
共鳴面 $2t+\rho^2-1=0$ の全てのH3候補を排除できる。
\end{proposition}
\begin{proof}
前節で残ったのは $\rho^2=-1,-3$ の二曲線だけである。
新しい入力
\[
 \texttt{cubic\_order2\_22\_secant\_main\_functionals\_input.json}
\]
は、三パラメータ $\rho,s,t$ の多項式係数を持つ八つの四次汎関数を与える。
各汎関数について、元の六二次関係を三位置と十六個の外側語へ挿入した
全288行を消す多項式恒等式を直接検証する。
関係の分母は $(\rho-1)(\rho+1)^2$ で払えるので、必要atlasの全域で有効である。
消去で選んだgeneric商基底の極を用いず、FLINTによる生成計算を再実行する必要もない。

これらから正自由三次32語への正常行列を作り、$\rho=i,t=1$ と
$\rho=\sqrt{-3},t=2$ に特殊化する。
係数はそれぞれ $\mathbb Q(i)[s]$、$\mathbb Q(\sqrt{-3})[s]$ に属する。
$s\ne0$ 上で各行の共通 $s$ 冪だけを除き、三つの8minorを取ると、そのgcdは単元を除いて
\[
 \begin{array}{c|c}
 \rho=i&s^3\\
 \rho=\sqrt{-3}&s^3\left(s^2+\frac{2-2\sqrt{-3}}3\right)
 \end{array}
\]
となる。従って残余は第一曲線の $s=0$ と第二曲線の
$s=0$、$3s^2+2-2\sqrt{-3}=0$ のみである。
第二曲線の $s=0$ は既に排除した $a=0$ 境界である。
第一曲線の $s=0$ を元の関係から $\mathbb Q(i)$ 上で再消去すると、
三次商次元20、四次商次元36となり、AS4の必要Hilbert値に反する。
第二曲線の残る点は
\[
 \rho=\sqrt{-3},\qquad s=\frac{1+\sqrt{-3}}{\sqrt3}
\]
を $\mathbb Q(\sqrt3,i)$ 上で調べればよい。
元の全関係の厳密消去により三次商次元20、四次商次元35、正常rank8を得る。
この数体の埋込みが $\rho$ の両共役と $s$ の両符号を全て覆う。
第一曲線の $\rho=-i$ も体の共役で覆われる。

AS4のHilbert値を満たす全残余点で正常核次元は2以下となり、既存の全値因子pencilが次元2を与える。
従って全正常核がそのpencilに等しく、全切断が一次因子を持つ。
前節の開集合coverと合わせて共鳴面全体を閉じた。

\texttt{verify\_cubic\_order2\_22\_secant\_resonance\_final.py}
は多項式汎関数の全元関係恒等式、両数体上のminorとgcd、有限点の元関係消去を再生する。
必要なのはSymPyの厳密演算であり、近似根や有限標本からの外挿ではない。
残る主chartは $a\ne0$、$2t+\rho^2-1\ne0$ である。
この主chart及び全H3の分類は未完成である。
\end{proof}
